

\begin{derivation}[从\eqnrefs{eq:maxwell-lippmann-schwinger-dp68}到\eqnrefs{eq:maxwell_scattered_field_t_operator}]
把背景 Maxwell Green 算符记为 $\mathbf G_0^R$，介质相对于背景的散射势记为 $\mathsf V_{\rm EM}$。正文的 Lippmann--Schwinger 方程是
\[
\widetilde{\mathbf E}
=\widetilde{\mathbf E}_{\rm inc}
+\mathbf G_0^R\mathsf V_{\rm EM}\widetilde{\mathbf E}.
\]
先不要展开级数，而把未知总场的项移到左边：
\[
(\mathbf I-\mathbf G_0^R\mathsf V_{\rm EM})\widetilde{\mathbf E}
=\widetilde{\mathbf E}_{\rm inc}.
\]
只要所考察复频率不位于该算符的谱上，就存在逆算符，故
\[
\widetilde{\mathbf E}
=(\mathbf I-\mathbf G_0^R\mathsf V_{\rm EM})^{-1}
\widetilde{\mathbf E}_{\rm inc}.
\]
散射场因此为
\begin{align*}
\widetilde{\mathbf E}_{\rm sc}
&\equiv\widetilde{\mathbf E}-\widetilde{\mathbf E}_{\rm inc}\\
&=(\mathbf I-\mathbf G_0^R\mathsf V_{\rm EM})^{-1}
\mathbf G_0^R\mathsf V_{\rm EM}\widetilde{\mathbf E}_{\rm inc}.
\end{align*}
现在用可直接验证的算符恒等式
\[
(\mathbf I-AB)^{-1}A=A(\mathbf I-BA)^{-1}.
\]
令 $A=\mathbf G_0^R$、$B=\mathsf V_{\rm EM}$，便有
\[
\widetilde{\mathbf E}_{\rm sc}
=\mathbf G_0^R
\mathsf V_{\rm EM}
(\mathbf I-\mathbf G_0^R\mathsf V_{\rm EM})^{-1}
\widetilde{\mathbf E}_{\rm inc}.
\]
因此定义
\[
\boxed{
\mathsf T_{\rm EM}
\equiv\mathsf V_{\rm EM}
(\mathbf I-\mathbf G_0^R\mathsf V_{\rm EM})^{-1}}
\]
后立即得到
\[
\boxed{
\widetilde{\mathbf E}_{\rm sc}
=\mathbf G_0^R\mathsf T_{\rm EM}\widetilde{\mathbf E}_{\rm inc}},
\]
即正文\eqnrefs{eq:maxwell_scattered_field_t_operator}。

多次散射级数只是这个精确逆算符在收敛区域内的 Neumann 展开。若某个算符范数满足 $\lVert\mathbf G_0^R\mathsf V_{\rm EM}\rVert<1$，则
\begin{align*}
(\mathbf I-\mathbf G_0^R\mathsf V_{\rm EM})^{-1}
&=\sum_{n=0}^{\infty}(\mathbf G_0^R\mathsf V_{\rm EM})^n,\\
\mathsf T_{\rm EM}
&=\mathsf V_{\rm EM}
+\mathsf V_{\rm EM}\mathbf G_0^R\mathsf V_{\rm EM}
+\mathsf V_{\rm EM}\mathbf G_0^R\mathsf V_{\rm EM}\mathbf G_0^R\mathsf V_{\rm EM}+\cdots.
\end{align*}
第 $n$ 个额外因子 $\mathbf G_0^R\mathsf V_{\rm EM}$ 对应“再传播一次、再散射一次”。即使 Neumann 级数本身不收敛，$\mathsf T_{\rm EM}$ 仍由前面的算符逆定义；其极点对应齐次方程
\[
(\mathbf I-\mathbf G_0^R\mathsf V_{\rm EM})\mathbf E=0
\]
出现非零解，也就是
$\ker(\mathbf I-\mathbf G_0^R\mathsf V_{\rm EM})\neq\{0\}$；在有限维离散化后，这才等价于相应行列式为零。它对应无外源条件下仍可存在的本征模或共振。这样 $T$ 算符同时统一了弱散射 Born 级数与强多重散射的谱结构。
\end{derivation}

\begin{derivation}[从\eqnrefs{eq:maxwell-bilinear-reciprocity-dp68}到\eqnrefs{eq:maxwell-green-reciprocity-r16}]

正文\eqnrefs{eq:maxwell-bilinear-reciprocity-dp68} 已经把介质互易性写成 Maxwell 算符的双线性对称关系。要把它变成 Green 张量的源点交换关系，只需让两组测试场分别由两个点源产生。

取任意常矢量 \(\mathbf a,\mathbf b\)，令
\begin{equation*}
\mathbf u(\mathbf s)=\mathbf G^R(\mathbf s,\mathbf r;\omega)\cdot\mathbf a,
\qquad
\mathbf v(\mathbf s)=\mathbf G^R(\mathbf s,\mathbf r';\omega)\cdot\mathbf b.
\end{equation*}
由 Green 方程，
\begin{equation*}
\mathsf L_{\rm EM}\mathbf u
=\mathbf a\,\delta^{(3)}(\mathbf s-\mathbf r),
\qquad
\mathsf L_{\rm EM}\mathbf v
=\mathbf b\,\delta^{(3)}(\mathbf s-\mathbf r').
\end{equation*}
代入正文\eqnrefs{eq:maxwell-bilinear-reciprocity-dp68} 后，两个空间积分分别被 \(\delta\) 函数取样：
\begin{equation*}
\mathbf a\cdot\mathbf G^R(\mathbf r,\mathbf r';\omega)\cdot\mathbf b
=\mathbf b\cdot\mathbf G^R(\mathbf r',\mathbf r;\omega)\cdot\mathbf a.
\end{equation*}
因为 \(\mathbf a,\mathbf b\) 任意，所以
\begin{equation*}
\mathbf G^R(\mathbf r,\mathbf r';\omega)
=[\mathbf G^R(\mathbf r',\mathbf r;\omega)]^T,
\end{equation*}
即正文\eqnrefs{eq:maxwell-green-reciprocity-r16}。
\end{derivation}

\begin{derivation}[从\eqnrefs{eq:GR-Maxwell-SI-r12}到\eqnrefs{eq:green-optical-identity-SI-r12}]
令固定实频 $\omega>0$ 下的 Maxwell 算符为 $\mathsf L_{\rm EM}$，并定义
\[
\mathbf G^R=\mathsf L_{\rm EM}^{-1},
\qquad
\mathbf G^{R\dagger}=(\mathsf L_{\rm EM}^{\dagger})^{-1}.
\]
对任意两个可逆算符 $A,B$，有两种等价的预解式恒等式
\[
A^{-1}-B^{-1}
=A^{-1}(B-A)B^{-1}
=B^{-1}(B-A)A^{-1}.
\]
第一式例如由
$A^{-1}BB^{-1}-A^{-1}AB^{-1}=A^{-1}-B^{-1}$ 逐项得到。取
$A=\mathsf L_{\rm EM}$、$B=\mathsf L_{\rm EM}^{\dagger}$，先用第一种次序得到
\begin{align*}
\mathbf G^R-\mathbf G^{R\dagger}
&=\mathbf G^R
(\mathsf L_{\rm EM}^{\dagger}-\mathsf L_{\rm EM})
\mathbf G^{R\dagger}.
\end{align*}
材料的吸收与色散信息集中在算符差 $\mathsf L_{\rm EM}^{\dagger}-\mathsf L_{\rm EM}$ 中。对局域电介质且磁响应无损的情形，
\[
\mathsf L_{\rm EM}
=\nabla\times\nabla\times
-\frac{\omega^2}{c^2}\boldsymbol\varepsilon_r,
\]
故
\[
\mathsf L_{\rm EM}^{\dagger}-\mathsf L_{\rm EM}
=\frac{\omega^2}{c^2}
(\boldsymbol\varepsilon_r-\boldsymbol\varepsilon_r^{\dagger})
=2\ii\frac{\omega^2}{c^2}
\operatorname{Im}\boldsymbol\varepsilon_r.
\]
代回并除以 $2\ii$：
\[
\operatorname{Im}\mathbf G^R
=\frac{\omega^2}{c^2}
\mathbf G^R
(\operatorname{Im}\boldsymbol\varepsilon_r)
\mathbf G^{R\dagger}.
\]
把算符乘积写成坐标核。若介电张量局域，
$\operatorname{Im}\boldsymbol\varepsilon_r$ 的核为
$\operatorname{Im}\boldsymbol\varepsilon_r(\mathbf s,\omega)
\delta^{(3)}(\mathbf s-\mathbf s')$，于是
\[
\boxed{
\operatorname{Im}\mathbf G^R(\mathbf r,\mathbf r';\omega)
=\frac{\omega^2}{c^2}\int\dd^3s\,
\mathbf G^R(\mathbf r,\mathbf s;\omega)
\cdot\operatorname{Im}\boldsymbol\varepsilon_r(\mathbf s,\omega)
\cdot\mathbf G^{R\dagger}(\mathbf r',\mathbf s;\omega)}.
\]
这就是正文\eqnrefs{eq:green-optical-identity-SI-r12}。

该恒等式还解释了为什么 $\operatorname{Im}\mathbf G^R$ 能作为损失谱权重。为把正定性写成最直观的场能量形式，改用上面第二种等价次序，亦有
\[
\operatorname{Im}\mathbf G^R
=\frac{\omega^2}{c^2}
\mathbf G^{R\dagger}
(\operatorname{Im}\boldsymbol\varepsilon_r)
\mathbf G^R.
\]
对任意测试电流 $\mathbf J$，定义由它激发的场
\[
\mathbf E_J(\mathbf s)
\equiv\int\dd^3r'\,\mathbf G^R(\mathbf s,\mathbf r';\omega)\cdot\mathbf J(\mathbf r').
\]
于是左右与 $\mathbf J^*,\mathbf J$ 收缩得到
\[
\iint\dd^3r\dd^3r'\,
\mathbf J^*(\mathbf r)\cdot\operatorname{Im}\mathbf G^R(\mathbf r,\mathbf r')\cdot\mathbf J(\mathbf r')
=\frac{\omega^2}{c^2}\int\dd^3s\,
\mathbf E_J^*(\mathbf s)\cdot
\operatorname{Im}\boldsymbol\varepsilon_r(\mathbf s)\cdot\mathbf E_J(\mathbf s).
\]
对被动介质，$\operatorname{Im}\boldsymbol\varepsilon_r\succeq0$，所以右边非负。因而 Green 张量的反厄米部分不是一个形式上的“取虚部”规则，而是由材料局域耗散经过两次传播映射回源空间后得到的正半定谱核；这正是后文 EELS、LDOS 与涨落--耗散关系共同使用的数学结构。
\end{derivation}
